\subsection*{Вариант 1 для  Рыльской Эвилины Ильгамовны. Дроби \\ \today   \currenttime} 
  \begin{enumerate}
\item Найди число, если его $\frac{1}{9}$  часть равна 54. Нарисуй схему.  \\ \nopagebreak \begin{tikzpicture}
 \draw[step=0.7cm ,gray,very thin] (0,0) grid (18,5);
 % Обводка внешнего контура  \draw[black, thick] (0,0) rectangle (18,5);
 \end{tikzpicture}\item Найди $\frac{1}{5}$ часть от 50. Нарисуй схему. \\ \nopagebreak \begin{tikzpicture}
 \draw[step=0.7cm ,gray,very thin] (0,0) grid (18,5);
 % Обводка внешнего контура  \draw[black, thick] (0,0) rectangle (18,5);
 \end{tikzpicture}\item Найди число, если его $\frac{5}{8}$  часть равна 40. Нарисуй схему.  \\ \nopagebreak \begin{tikzpicture}
 \draw[step=0.7cm ,gray,very thin] (0,0) grid (18,5);
 % Обводка внешнего контура  \draw[black, thick] (0,0) rectangle (18,5);
 \end{tikzpicture}\item Найди $\frac{5}{15}$ часть от 75.  Нарисуй схему. \\ \nopagebreak \begin{tikzpicture}
 \draw[step=0.7cm ,gray,very thin] (0,0) grid (18,5);
 % Обводка внешнего контура  \draw[black, thick] (0,0) rectangle (18,5);
 \end{tikzpicture}\item Найди число, если его  12 \%  равны 10800. Нарисуй схему.  \\ \nopagebreak \begin{tikzpicture}
 \draw[step=0.7cm ,gray,very thin] (0,0) grid (18,5);
 % Обводка внешнего контура  \draw[black, thick] (0,0) rectangle (18,5);
 \end{tikzpicture}\item Найди 8 \%  от 12800.  Нарисуй схему. \\ \nopagebreak \begin{tikzpicture}
 \draw[step=0.7cm ,gray,very thin] (0,0) grid (18,5);
 % Обводка внешнего контура  \draw[black, thick] (0,0) rectangle (18,5);
 \end{tikzpicture}\end{enumerate}  
\newpage
\subsection*{Вариант 2 для  Абелюк Марии. Дроби \\ \today   \currenttime} 
  \begin{enumerate}
\item Найди число, если его $\frac{1}{4}$  часть равна 20. Нарисуй схему.  \\ \nopagebreak \begin{tikzpicture}
 \draw[step=0.7cm ,gray,very thin] (0,0) grid (18,5);
 % Обводка внешнего контура  \draw[black, thick] (0,0) rectangle (18,5);
 \end{tikzpicture}\item Найди $\frac{1}{14}$ часть от 140. Нарисуй схему. \\ \nopagebreak \begin{tikzpicture}
 \draw[step=0.7cm ,gray,very thin] (0,0) grid (18,5);
 % Обводка внешнего контура  \draw[black, thick] (0,0) rectangle (18,5);
 \end{tikzpicture}\item Найди число, если его $\frac{5}{16}$  часть равна 80. Нарисуй схему.  \\ \nopagebreak \begin{tikzpicture}
 \draw[step=0.7cm ,gray,very thin] (0,0) grid (18,5);
 % Обводка внешнего контура  \draw[black, thick] (0,0) rectangle (18,5);
 \end{tikzpicture}\item Найди $\frac{8}{14}$ часть от 112.  Нарисуй схему. \\ \nopagebreak \begin{tikzpicture}
 \draw[step=0.7cm ,gray,very thin] (0,0) grid (18,5);
 % Обводка внешнего контура  \draw[black, thick] (0,0) rectangle (18,5);
 \end{tikzpicture}\item Найди число, если его  20 \%  равны 14000. Нарисуй схему.  \\ \nopagebreak \begin{tikzpicture}
 \draw[step=0.7cm ,gray,very thin] (0,0) grid (18,5);
 % Обводка внешнего контура  \draw[black, thick] (0,0) rectangle (18,5);
 \end{tikzpicture}\item Найди 6 \%  от 7800.  Нарисуй схему. \\ \nopagebreak \begin{tikzpicture}
 \draw[step=0.7cm ,gray,very thin] (0,0) grid (18,5);
 % Обводка внешнего контура  \draw[black, thick] (0,0) rectangle (18,5);
 \end{tikzpicture}\end{enumerate}  
\newpage
\subsection*{Вариант 3 для  Гондарева Ильи. Дроби \\ \today   \currenttime} 
  \begin{enumerate}
\item Найди число, если его $\frac{1}{10}$  часть равна 60. Нарисуй схему.  \\ \nopagebreak \begin{tikzpicture}
 \draw[step=0.7cm ,gray,very thin] (0,0) grid (18,5);
 % Обводка внешнего контура  \draw[black, thick] (0,0) rectangle (18,5);
 \end{tikzpicture}\item Найди $\frac{1}{8}$ часть от 104. Нарисуй схему. \\ \nopagebreak \begin{tikzpicture}
 \draw[step=0.7cm ,gray,very thin] (0,0) grid (18,5);
 % Обводка внешнего контура  \draw[black, thick] (0,0) rectangle (18,5);
 \end{tikzpicture}\item Найди число, если его $\frac{10}{22}$  часть равна 220. Нарисуй схему.  \\ \nopagebreak \begin{tikzpicture}
 \draw[step=0.7cm ,gray,very thin] (0,0) grid (18,5);
 % Обводка внешнего контура  \draw[black, thick] (0,0) rectangle (18,5);
 \end{tikzpicture}\item Найди $\frac{7}{14}$ часть от 98.  Нарисуй схему. \\ \nopagebreak \begin{tikzpicture}
 \draw[step=0.7cm ,gray,very thin] (0,0) grid (18,5);
 % Обводка внешнего контура  \draw[black, thick] (0,0) rectangle (18,5);
 \end{tikzpicture}\item Найди число, если его  3 \%  равны 3900. Нарисуй схему.  \\ \nopagebreak \begin{tikzpicture}
 \draw[step=0.7cm ,gray,very thin] (0,0) grid (18,5);
 % Обводка внешнего контура  \draw[black, thick] (0,0) rectangle (18,5);
 \end{tikzpicture}\item Найди 3 \%  от 6600.  Нарисуй схему. \\ \nopagebreak \begin{tikzpicture}
 \draw[step=0.7cm ,gray,very thin] (0,0) grid (18,5);
 % Обводка внешнего контура  \draw[black, thick] (0,0) rectangle (18,5);
 \end{tikzpicture}\end{enumerate}  
\newpage
\subsection*{Вариант 4 для  Аношина Матвея Антоновича. Дроби \\ \today   \currenttime} 
  \begin{enumerate}
\item Найди число, если его $\frac{1}{6}$  часть равна 54. Нарисуй схему.  \\ \nopagebreak \begin{tikzpicture}
 \draw[step=0.7cm ,gray,very thin] (0,0) grid (18,5);
 % Обводка внешнего контура  \draw[black, thick] (0,0) rectangle (18,5);
 \end{tikzpicture}\item Найди $\frac{1}{2}$ часть от 18. Нарисуй схему. \\ \nopagebreak \begin{tikzpicture}
 \draw[step=0.7cm ,gray,very thin] (0,0) grid (18,5);
 % Обводка внешнего контура  \draw[black, thick] (0,0) rectangle (18,5);
 \end{tikzpicture}\item Найди число, если его $\frac{4}{12}$  часть равна 48. Нарисуй схему.  \\ \nopagebreak \begin{tikzpicture}
 \draw[step=0.7cm ,gray,very thin] (0,0) grid (18,5);
 % Обводка внешнего контура  \draw[black, thick] (0,0) rectangle (18,5);
 \end{tikzpicture}\item Найди $\frac{6}{14}$ часть от 84.  Нарисуй схему. \\ \nopagebreak \begin{tikzpicture}
 \draw[step=0.7cm ,gray,very thin] (0,0) grid (18,5);
 % Обводка внешнего контура  \draw[black, thick] (0,0) rectangle (18,5);
 \end{tikzpicture}\item Найди число, если его  11 \%  равны 3300. Нарисуй схему.  \\ \nopagebreak \begin{tikzpicture}
 \draw[step=0.7cm ,gray,very thin] (0,0) grid (18,5);
 % Обводка внешнего контура  \draw[black, thick] (0,0) rectangle (18,5);
 \end{tikzpicture}\item Найди 10 \%  от 22000.  Нарисуй схему. \\ \nopagebreak \begin{tikzpicture}
 \draw[step=0.7cm ,gray,very thin] (0,0) grid (18,5);
 % Обводка внешнего контура  \draw[black, thick] (0,0) rectangle (18,5);
 \end{tikzpicture}\end{enumerate}  
\newpage
\subsection*{Вариант 5 для  Егоровой Мирославы Даниловны. Дроби \\ \today   \currenttime} 
  \begin{enumerate}
\item Найди число, если его $\frac{1}{10}$  часть равна 90. Нарисуй схему.  \\ \nopagebreak \begin{tikzpicture}
 \draw[step=0.7cm ,gray,very thin] (0,0) grid (18,5);
 % Обводка внешнего контура  \draw[black, thick] (0,0) rectangle (18,5);
 \end{tikzpicture}\item Найди $\frac{1}{14}$ часть от 112. Нарисуй схему. \\ \nopagebreak \begin{tikzpicture}
 \draw[step=0.7cm ,gray,very thin] (0,0) grid (18,5);
 % Обводка внешнего контура  \draw[black, thick] (0,0) rectangle (18,5);
 \end{tikzpicture}\item Найди число, если его $\frac{2}{6}$  часть равна 12. Нарисуй схему.  \\ \nopagebreak \begin{tikzpicture}
 \draw[step=0.7cm ,gray,very thin] (0,0) grid (18,5);
 % Обводка внешнего контура  \draw[black, thick] (0,0) rectangle (18,5);
 \end{tikzpicture}\item Найди $\frac{12}{20}$ часть от 240.  Нарисуй схему. \\ \nopagebreak \begin{tikzpicture}
 \draw[step=0.7cm ,gray,very thin] (0,0) grid (18,5);
 % Обводка внешнего контура  \draw[black, thick] (0,0) rectangle (18,5);
 \end{tikzpicture}\item Найди число, если его  16 \%  равны 16000. Нарисуй схему.  \\ \nopagebreak \begin{tikzpicture}
 \draw[step=0.7cm ,gray,very thin] (0,0) grid (18,5);
 % Обводка внешнего контура  \draw[black, thick] (0,0) rectangle (18,5);
 \end{tikzpicture}\item Найди 6 \%  от 12600.  Нарисуй схему. \\ \nopagebreak \begin{tikzpicture}
 \draw[step=0.7cm ,gray,very thin] (0,0) grid (18,5);
 % Обводка внешнего контура  \draw[black, thick] (0,0) rectangle (18,5);
 \end{tikzpicture}\end{enumerate}  
\newpage
\subsection*{Вариант 6 для  Карнауха Тимофея Павловича. Дроби \\ \today   \currenttime} 
  \begin{enumerate}
\item Найди число, если его $\frac{1}{9}$  часть равна 54. Нарисуй схему.  \\ \nopagebreak \begin{tikzpicture}
 \draw[step=0.7cm ,gray,very thin] (0,0) grid (18,5);
 % Обводка внешнего контура  \draw[black, thick] (0,0) rectangle (18,5);
 \end{tikzpicture}\item Найди $\frac{1}{12}$ часть от 144. Нарисуй схему. \\ \nopagebreak \begin{tikzpicture}
 \draw[step=0.7cm ,gray,very thin] (0,0) grid (18,5);
 % Обводка внешнего контура  \draw[black, thick] (0,0) rectangle (18,5);
 \end{tikzpicture}\item Найди число, если его $\frac{9}{13}$  часть равна 117. Нарисуй схему.  \\ \nopagebreak \begin{tikzpicture}
 \draw[step=0.7cm ,gray,very thin] (0,0) grid (18,5);
 % Обводка внешнего контура  \draw[black, thick] (0,0) rectangle (18,5);
 \end{tikzpicture}\item Найди $\frac{7}{14}$ часть от 98.  Нарисуй схему. \\ \nopagebreak \begin{tikzpicture}
 \draw[step=0.7cm ,gray,very thin] (0,0) grid (18,5);
 % Обводка внешнего контура  \draw[black, thick] (0,0) rectangle (18,5);
 \end{tikzpicture}\item Найди число, если его  10 \%  равны 21000. Нарисуй схему.  \\ \nopagebreak \begin{tikzpicture}
 \draw[step=0.7cm ,gray,very thin] (0,0) grid (18,5);
 % Обводка внешнего контура  \draw[black, thick] (0,0) rectangle (18,5);
 \end{tikzpicture}\item Найди 10 \%  от 6000.  Нарисуй схему. \\ \nopagebreak \begin{tikzpicture}
 \draw[step=0.7cm ,gray,very thin] (0,0) grid (18,5);
 % Обводка внешнего контура  \draw[black, thick] (0,0) rectangle (18,5);
 \end{tikzpicture}\end{enumerate}  
\newpage
\subsection*{Вариант 7 для  Костенко Каролины Романовны. Дроби \\ \today   \currenttime} 
  \begin{enumerate}
\item Найди число, если его $\frac{1}{8}$  часть равна 40. Нарисуй схему.  \\ \nopagebreak \begin{tikzpicture}
 \draw[step=0.7cm ,gray,very thin] (0,0) grid (18,5);
 % Обводка внешнего контура  \draw[black, thick] (0,0) rectangle (18,5);
 \end{tikzpicture}\item Найди $\frac{1}{11}$ часть от 44. Нарисуй схему. \\ \nopagebreak \begin{tikzpicture}
 \draw[step=0.7cm ,gray,very thin] (0,0) grid (18,5);
 % Обводка внешнего контура  \draw[black, thick] (0,0) rectangle (18,5);
 \end{tikzpicture}\item Найди число, если его $\frac{10}{19}$  часть равна 190. Нарисуй схему.  \\ \nopagebreak \begin{tikzpicture}
 \draw[step=0.7cm ,gray,very thin] (0,0) grid (18,5);
 % Обводка внешнего контура  \draw[black, thick] (0,0) rectangle (18,5);
 \end{tikzpicture}\item Найди $\frac{6}{11}$ часть от 66.  Нарисуй схему. \\ \nopagebreak \begin{tikzpicture}
 \draw[step=0.7cm ,gray,very thin] (0,0) grid (18,5);
 % Обводка внешнего контура  \draw[black, thick] (0,0) rectangle (18,5);
 \end{tikzpicture}\item Найди число, если его  13 \%  равны 16900. Нарисуй схему.  \\ \nopagebreak \begin{tikzpicture}
 \draw[step=0.7cm ,gray,very thin] (0,0) grid (18,5);
 % Обводка внешнего контура  \draw[black, thick] (0,0) rectangle (18,5);
 \end{tikzpicture}\item Найди 5 \%  от 5000.  Нарисуй схему. \\ \nopagebreak \begin{tikzpicture}
 \draw[step=0.7cm ,gray,very thin] (0,0) grid (18,5);
 % Обводка внешнего контура  \draw[black, thick] (0,0) rectangle (18,5);
 \end{tikzpicture}\end{enumerate}  
\newpage
\subsection*{Вариант 8 для  Минца Марка Евгеньевича. Дроби \\ \today   \currenttime} 
  \begin{enumerate}
\item Найди число, если его $\frac{1}{5}$  часть равна 35. Нарисуй схему.  \\ \nopagebreak \begin{tikzpicture}
 \draw[step=0.7cm ,gray,very thin] (0,0) grid (18,5);
 % Обводка внешнего контура  \draw[black, thick] (0,0) rectangle (18,5);
 \end{tikzpicture}\item Найди $\frac{1}{4}$ часть от 36. Нарисуй схему. \\ \nopagebreak \begin{tikzpicture}
 \draw[step=0.7cm ,gray,very thin] (0,0) grid (18,5);
 % Обводка внешнего контура  \draw[black, thick] (0,0) rectangle (18,5);
 \end{tikzpicture}\item Найди число, если его $\frac{9}{18}$  часть равна 162. Нарисуй схему.  \\ \nopagebreak \begin{tikzpicture}
 \draw[step=0.7cm ,gray,very thin] (0,0) grid (18,5);
 % Обводка внешнего контура  \draw[black, thick] (0,0) rectangle (18,5);
 \end{tikzpicture}\item Найди $\frac{6}{12}$ часть от 72.  Нарисуй схему. \\ \nopagebreak \begin{tikzpicture}
 \draw[step=0.7cm ,gray,very thin] (0,0) grid (18,5);
 % Обводка внешнего контура  \draw[black, thick] (0,0) rectangle (18,5);
 \end{tikzpicture}\item Найди число, если его  16 \%  равны 25600. Нарисуй схему.  \\ \nopagebreak \begin{tikzpicture}
 \draw[step=0.7cm ,gray,very thin] (0,0) grid (18,5);
 % Обводка внешнего контура  \draw[black, thick] (0,0) rectangle (18,5);
 \end{tikzpicture}\item Найди 5 \%  от 2000.  Нарисуй схему. \\ \nopagebreak \begin{tikzpicture}
 \draw[step=0.7cm ,gray,very thin] (0,0) grid (18,5);
 % Обводка внешнего контура  \draw[black, thick] (0,0) rectangle (18,5);
 \end{tikzpicture}\end{enumerate}  
\newpage
\subsection*{Вариант 9 для  Саргсян Эмилии Арамовны. Дроби \\ \today   \currenttime} 
  \begin{enumerate}
\item Найди число, если его $\frac{1}{10}$  часть равна 60. Нарисуй схему.  \\ \nopagebreak \begin{tikzpicture}
 \draw[step=0.7cm ,gray,very thin] (0,0) grid (18,5);
 % Обводка внешнего контура  \draw[black, thick] (0,0) rectangle (18,5);
 \end{tikzpicture}\item Найди $\frac{1}{8}$ часть от 96. Нарисуй схему. \\ \nopagebreak \begin{tikzpicture}
 \draw[step=0.7cm ,gray,very thin] (0,0) grid (18,5);
 % Обводка внешнего контура  \draw[black, thick] (0,0) rectangle (18,5);
 \end{tikzpicture}\item Найди число, если его $\frac{9}{20}$  часть равна 180. Нарисуй схему.  \\ \nopagebreak \begin{tikzpicture}
 \draw[step=0.7cm ,gray,very thin] (0,0) grid (18,5);
 % Обводка внешнего контура  \draw[black, thick] (0,0) rectangle (18,5);
 \end{tikzpicture}\item Найди $\frac{12}{22}$ часть от 264.  Нарисуй схему. \\ \nopagebreak \begin{tikzpicture}
 \draw[step=0.7cm ,gray,very thin] (0,0) grid (18,5);
 % Обводка внешнего контура  \draw[black, thick] (0,0) rectangle (18,5);
 \end{tikzpicture}\item Найди число, если его  18 \%  равны 32400. Нарисуй схему.  \\ \nopagebreak \begin{tikzpicture}
 \draw[step=0.7cm ,gray,very thin] (0,0) grid (18,5);
 % Обводка внешнего контура  \draw[black, thick] (0,0) rectangle (18,5);
 \end{tikzpicture}\item Найди 12 \%  от 19200.  Нарисуй схему. \\ \nopagebreak \begin{tikzpicture}
 \draw[step=0.7cm ,gray,very thin] (0,0) grid (18,5);
 % Обводка внешнего контура  \draw[black, thick] (0,0) rectangle (18,5);
 \end{tikzpicture}\end{enumerate}  
\newpage
\subsection*{Вариант 10 для  Строкатовой Киры Павловны. Дроби \\ \today   \currenttime} 
  \begin{enumerate}
\item Найди число, если его $\frac{1}{11}$  часть равна 77. Нарисуй схему.  \\ \nopagebreak \begin{tikzpicture}
 \draw[step=0.7cm ,gray,very thin] (0,0) grid (18,5);
 % Обводка внешнего контура  \draw[black, thick] (0,0) rectangle (18,5);
 \end{tikzpicture}\item Найди $\frac{1}{2}$ часть от 14. Нарисуй схему. \\ \nopagebreak \begin{tikzpicture}
 \draw[step=0.7cm ,gray,very thin] (0,0) grid (18,5);
 % Обводка внешнего контура  \draw[black, thick] (0,0) rectangle (18,5);
 \end{tikzpicture}\item Найди число, если его $\frac{5}{13}$  часть равна 65. Нарисуй схему.  \\ \nopagebreak \begin{tikzpicture}
 \draw[step=0.7cm ,gray,very thin] (0,0) grid (18,5);
 % Обводка внешнего контура  \draw[black, thick] (0,0) rectangle (18,5);
 \end{tikzpicture}\item Найди $\frac{4}{11}$ часть от 44.  Нарисуй схему. \\ \nopagebreak \begin{tikzpicture}
 \draw[step=0.7cm ,gray,very thin] (0,0) grid (18,5);
 % Обводка внешнего контура  \draw[black, thick] (0,0) rectangle (18,5);
 \end{tikzpicture}\item Найди число, если его  13 \%  равны 15600. Нарисуй схему.  \\ \nopagebreak \begin{tikzpicture}
 \draw[step=0.7cm ,gray,very thin] (0,0) grid (18,5);
 % Обводка внешнего контура  \draw[black, thick] (0,0) rectangle (18,5);
 \end{tikzpicture}\item Найди 10 \%  от 22000.  Нарисуй схему. \\ \nopagebreak \begin{tikzpicture}
 \draw[step=0.7cm ,gray,very thin] (0,0) grid (18,5);
 % Обводка внешнего контура  \draw[black, thick] (0,0) rectangle (18,5);
 \end{tikzpicture}\end{enumerate}  
\newpage
\subsection*{Вариант 11 для  Тавриной Марии Ивановны. Дроби \\ \today   \currenttime} 
  \begin{enumerate}
\item Найди число, если его $\frac{1}{11}$  часть равна 22. Нарисуй схему.  \\ \nopagebreak \begin{tikzpicture}
 \draw[step=0.7cm ,gray,very thin] (0,0) grid (18,5);
 % Обводка внешнего контура  \draw[black, thick] (0,0) rectangle (18,5);
 \end{tikzpicture}\item Найди $\frac{1}{2}$ часть от 8. Нарисуй схему. \\ \nopagebreak \begin{tikzpicture}
 \draw[step=0.7cm ,gray,very thin] (0,0) grid (18,5);
 % Обводка внешнего контура  \draw[black, thick] (0,0) rectangle (18,5);
 \end{tikzpicture}\item Найди число, если его $\frac{9}{14}$  часть равна 126. Нарисуй схему.  \\ \nopagebreak \begin{tikzpicture}
 \draw[step=0.7cm ,gray,very thin] (0,0) grid (18,5);
 % Обводка внешнего контура  \draw[black, thick] (0,0) rectangle (18,5);
 \end{tikzpicture}\item Найди $\frac{13}{20}$ часть от 260.  Нарисуй схему. \\ \nopagebreak \begin{tikzpicture}
 \draw[step=0.7cm ,gray,very thin] (0,0) grid (18,5);
 % Обводка внешнего контура  \draw[black, thick] (0,0) rectangle (18,5);
 \end{tikzpicture}\item Найди число, если его  9 \%  равны 11700. Нарисуй схему.  \\ \nopagebreak \begin{tikzpicture}
 \draw[step=0.7cm ,gray,very thin] (0,0) grid (18,5);
 % Обводка внешнего контура  \draw[black, thick] (0,0) rectangle (18,5);
 \end{tikzpicture}\item Найди 3 \%  от 4500.  Нарисуй схему. \\ \nopagebreak \begin{tikzpicture}
 \draw[step=0.7cm ,gray,very thin] (0,0) grid (18,5);
 % Обводка внешнего контура  \draw[black, thick] (0,0) rectangle (18,5);
 \end{tikzpicture}\end{enumerate}  
\newpage
